% tutorial from  https://texblog.org/2012/04/25/writing-a-cv-in-latex/

\documentclass[10pt]{article}

\usepackage[margin=3cm]{geometry}

\usepackage{hyperref}

\usepackage{array, xcolor}
\definecolor{lightgray}{gray}{0.8}
\newcolumntype{L}{>{\raggedleft}p{0.14\textwidth}}
\newcolumntype{R}{p{0.8\textwidth}}
\newcommand\VRule{\color{lightgray}\vrule width 0.5pt}

\usepackage{lipsum}

\pagestyle{empty}

\newcommand{\info}[2]{
  \noindent
  \parbox{8em}{\textbf{#1}} #2 \par}

\newcommand{\work}[4]{
\textnormal{#2}&\textbf{#1}\\
&\textnormal{#3}\\
&\textnormal{#4}\\
}

\newcommand{\education}[2]{
#1&#2\\
}

\newcommand{\skill}[2]{
#1&#2\\
&\\
}

\newcommand{\lang}[2]{
#1&#2\\
}

\begin{document}
\part*{Natan Camargos}

\info{Email}{\href{mailto:natanscamargos@gmail.com}{\texttt{natanscamargos@gmail.com}} }
\info{Celular}{\href{tel:+5585998056302}{+55 85 99805-6302}}
\info{Endereço}{Rua Aldemir Jucá, 44, 60824-100, Fortaleza/Ceará}
\info{Linkedin}{\href{www.linkedin.com/in/natan-camargos-7b8153122}{www.linkedin.com/in/natan-camargos-7b8153122}}
\info{Github}{\href{https://github.com/natancamargo}{https://github.com/natancamargo}}
\info{Stackoverflow}{\href{https://stackoverflow.com/users/6489712/natan-camargos}{https://stackoverflow.com/users/6489712/natan-camargos}}

\section*{Experiência}
\begin{tabular}{L!{\VRule}R}
\work
{Beyond Soluções}
{2022--2023}
{Desenvolvedor Senior}
{React.js $|$ Typescript $|$ Styled Components $|$ Storybook $|$ Jest $|$ RTL $|$ Cypress $|$
Webpack\newline
Desenvolvimento e manutenção de bibliotecas e ferramentas de frontend para
Português melhorar a experiência para o usuário final.
}
&\\
\work
{Fretebras}
{2020--2021}
{Desenvolvedor Senior}
{React.js $|$ Styled Components $|$ Storybook $|$ Jest $|$ RTL $|$ Cypress $|$ Webpack\newline
Frontend na plataforma de risco da FreteBras, que oferece validação de
motoristas através de OCR e consultas aos órgãos de trânsito brasileiros,
de modo a estimar o risco associado com cada frete. Desenvolvimento de
projetos cross-squads (e.g: Design System Components and CLI).
}
&\\
\work
{Hathor Network}
{2020--2021}
{Engenheiro de Software}
{React.js $|$ React Native $|$ Electron $|$ Node.js $|$ Jest\newline
Contribuição no desenvolvimento das carteiras mobile e desktop, além de
participação na construção de APIs e de bibliotecas no ecossistema Hathor.
}
&\\
\work
{Bitcoin.com}
{2020--2021}
{Engenheiro de Software}
{React.js $|$ Node.js $|$ Bitcoin Cash\newline
Desenvolvimento de uma carteira de ativos para a blockchain do Bitcoin
Cash. A carteira "Mint" foi criada durante um hackathon e atraiu bastante
o interesse da comunidade. A carteira foi divulgada em vários artigos,
incluindo os seguintes:
\begin{itemize}
  \item \href{https://news.bitcoin.com/how-to-create-custom-slp-tokens-with-the-bitcoin-com-mint/}
{https://news.bitcoin.com/how-to-create-custom-slp-tokens-with-the-bitcoin-com-mint/}
  \item \href{https://www.businesswire.com/news/home/20200213005745/en/Bitcoin.com-Launches-Mint-Token-Creation-
Platform-Set}{https://www.businesswire.com/news/home/20200213005745/en/Bitcoin.com-Launches-Mint-Token-Creation-
Platform-Set}
\end{itemize}
}
\end{tabular}

\section*{Educação}
\begin{tabular}{L!{\VRule}R}
\education{2009--2013}{
Universidade Federal do Ceará\newline
Engenharia de Teleinformática - Enfâse em Engenharia de Computação}
\end{tabular}

\section*{Habilidades}
\begin{tabular}{L!{\VRule}R}
\skill
{Linguagens de programação}
{Javascript, C/C++, Elisp, Java}
\skill
{Outros}
{CSS, HTML, Node, Ionic, Android, Markdown, Git, Emacs, Ubuntu}
\end{tabular}

\section*{Linguagens}
\begin{tabular}{L!{\VRule}R}
\lang{Português}{Nativo}
\lang{Inglês}{B2}
\end{tabular}
\end{document}
